% Hoja de vida Rodrigo Taborda -- RESUMEN de dos páginas (moderncv classic)
% Codificación: UTF-8. Versión completa: cv-rtaborda-espanol.tex
\documentclass[10pt,a4paper,sans]{moderncv}
\moderncvcolor{burgundy}
\moderncvstyle{classic}
\usepackage[utf8]{inputenc}
\usepackage[T1]{fontenc}
\usepackage[spanish]{babel}
\usepackage[scale=0.86]{geometry}
\setlength{\hintscolumnwidth}{2.6cm}

% ---------------------------------------------------------------------
% Datos personales
% ---------------------------------------------------------------------
    \name{Rodrigo}{Taborda}
    \title{Profesor asociado · Facultad de Administración}
    \address{Calle 21 \# 1-20, Facultad de Administración}{Universidad de los Andes, Bogotá, Colombia}
    \email{r.taborda@uniandes.edu.co}
    \homepage{www.rodrigotaborda.com}
    \extrainfo{Nacionalidad: colombiana}

% Entrada de publicación: {año}{título}{coautores}{fuente}
\newcommand{\pub}[4]{\cvitem{#1}{\textbf{#2}\ifx&#3&\else\ (con #3)\fi. #4}}

\begin{document}
\pagestyle{fancy}\fancyhf{}\renewcommand{\headrulewidth}{0pt}
\fancyfoot[C]{\addressfont\color{color2}Rodrigo Taborda --- Hoja de vida (resumen) --- \thepage}
\makecvtitle

    \section{Cargos}
        \cventry{2018--}{Profesor asociado}{Facultad de Administración, Universidad de los Andes}{}{}{}
        \cventry{2013--2018}{Profesor asistente}{Facultad de Administración, Universidad de los Andes}{}{}{}
        \cventry{2010--2013}{Profesor}{Facultad de Economía, Universidad del Rosario}{}{}{}
        \cventry{2010--2023}{Editor}{\emph{Revista de Economía del Rosario}}{}{}{}
        \cventry{2008--2010}{Consultor}{ACIL-Tasman, Canberra, Australia}{}{}{}
        \cventry{2002--2004}{Profesor}{Universidad del Rosario}{}{}{}

    \section{Educación}
        \cventry{2009}{Ph.D. Economics}{Australian National University}{}{}{Tesis: \emph{Manufacturing industry performance after structural reforms. The case of Colombia}.}
        \cventry{2005}{M.A. Economics}{Australian National University}{}{}{}
        \cventry{2004}{Maestría en Economía}{Universidad del Rosario}{}{}{}
        \cventry{2000}{Economía}{Universidad de Antioquia}{}{}{}

    \section{Publicaciones}
        \subsection{Artículos en revistas académicas}
            \pub{2027}{A test of auditor performance under the service paradigm: do emotional intelligence and national culture add to our understanding?}{Rosa González, Stephen B. Salter y Felipe Janica}{\emph{Journal of International Accounting, Auditing and Taxation}, en prensa}
            \pub{2025}{Time poverty: an unintended consequence of women participation in farmers' associations}{Andrés Barrios y Ximena Rueda}{\emph{Social Indicators Research}, 178(1), 225--253}
            \pub{2023}{Fund flows, stock markets, and economic policy uncertainty: from the perspective of a CIVETS nation}{Joseph J. French y Philipp D. Schaberl}{\emph{Emerging Markets Finance and Trade}, 60(5), 967--982}
            \pub{2022}{Digitalization and e-government in the lives of urban migrants: evidence from Bogota}{Charles Martin, Sonia Camacho y Constantine Ruhe}{\emph{Policy \& Internet}, 14(2)}
            \pub{2022}{Savings and gender in Colombia: ANOVA, simple linear regression, and extreme observations}{Luz Orozco}{\emph{SAGE Research Methods Cases}}
            \pub{2019}{How does the training of subsistence entrepreneurs affect their hope levels and perceived goal-attainment?}{Andrés Barrios y Ezequiel Reficco}{\emph{Education + Training}, 61(7/8)}
            \pub{2019}{Family firms and financial performance. A meta-regression analysis}{Maximiliano González y Juan Idrobo}{\emph{Academia Revista Latinoamericana de Administración}, 32(3)}
            \pub{2018}{Disentangling the relationship between liquidity and returns in Latin America}{Joseph French}{\emph{Global Finance Journal}, 36}
            \pub{2017}{Stock liquidity and second blockholder as drivers of corporate value: evidence from Latin America}{Carlos Pombo}{\emph{International Review of Economics and Finance}, 51}
            \pub{2015}{Innovation and productivity in the Colombian service and manufacturing industries}{Juan Gallego y Luis Gutiérrez}{\emph{Emerging Markets Finance and Trade}, 51}
            \pub{2015}{Procedural transparency in Latin American central banks under inflation targeting schemes. A text analysis of the minutes of the boards of directors}{}{\emph{Ensayos sobre Política Económica}, Banco de la República, 33}
            \pub{2014}{Insurance use and economic growth in Latin America. Panel data evidence}{Ángela Concha}{\emph{Lecturas de Economía}, 81, julio--diciembre}
            \pub{2013}{Bias in economic news: the reporting of the nominal exchange rate behavior in Colombia}{}{\emph{Economía -- Journal of LACEA}, 14(1)}
            \pub{2011}{Manufacturing employment and wage differentials after structural adjustment reforms in Colombia}{Juan Guataquí}{\emph{Lecturas de Economía}, 48, enero--junio}
            \pub{2008}{Ownership and control in Colombian corporations}{Carlos Pombo y Luis Gutiérrez}{\emph{Quarterly Review of Economics and Finance}, 48(1)}
            \pub{2006}{Technical instruments and conceptual definitions: theoretical and empirical implications of the change on the definition of unemployment in Colombia}{Juan Guataquí}{\emph{Revista de Economía del Rosario}, 9(1)}
            \pub{2006}{Performance, efficiency and productivity in Colombia's power distribution system: effects of the 1994 regulatory reform}{Carlos Pombo}{\emph{Energy Economics}, 28(3)}
        \subsection{Libros}
            \cvitem{}{\textbf{Ejercicios de econometría} (2016) y \textbf{Curso de macroeconomía} (2005), Editorial Universidad del Rosario.}

    \section{Docencia}
        \cvlistitem{29 cursos distintos dictados entre 2002 y 2026 en cuatro instituciones: 19 en la Universidad de los Andes y 8 en la Universidad del Rosario (Colombia), y 2 fuera de Colombia (IESA, Venezuela; Universidad Americana, Paraguay), además de tutorías en la Australian National University.}
        \cvlistitem{14 de pregrado, 14 de maestría y MBA y 1 de doctorado, más 2 programas de educación ejecutiva.}
        \cvlistitem{10 cursos de diseño original y 4 en línea (Maestría en Analítica y Gestión Financiera y MBA virtual).}
        \cvlistitem{Al menos 133 cursos-semestre en total, de los cuales 16 intersemestrales intensivos.}
        \cvlistitem{Docencia en español e inglés (Maestría Internacional en Finanzas con Nova SBE) y en codocencia.}

    \section{Dirección de tesis}
        \cvitem{}{Director de 2 tesis doctorales en curso (Doctorado en Administración), 3 de maestría (International Master in Finance NOVA--Uniandes; Maestría en Economía) y 2 de pregrado. Jurado o lector de 6 trabajos de grado de pregrado y 23 de maestría.}

    \section{Traducciones}
        \cvitem{}{18 libros traducidos del inglés al español (Banco Mundial, BID, Universidad del Rosario, Alfaomega, Mayol), entre ellos \emph{Oportunismo y buena voluntad} (Universidad del Rosario, en prensa), \emph{¿Quién gobierna a los excluidos?}, \emph{Género y ciudadanía en justicia transicional} (2024) y \emph{Las reformas agrarias en Latinoamérica} (2021); y 3 revisiones técnicas de traducciones (Banco Mundial, BID).}

    \section{Evaluador}
        \cvitem{Revistas}{Academia Revista Latinoamericana de Administración, Economic Change and Restructuring, Economic Papers, Economic Systems, Emerging Markets Finance and Trade, Empirical Economics, Energy Economics, Energy Journal, Ensayos sobre Política Económica, Estudios Gerenciales, International Review of Economics and Finance, Journal of Business Research, Lecturas de Economía, Managerial Finance, Revista de Economía del Rosario, Sociedad y Economía, Stata Journal.}
        \cvitem{Proyectos}{Swiss National Science Foundation (SNSF), COLCIENCIAS, Banco de la República (FPIT) y universidades colombianas. Jurado del Premio Otto de Greiff (2015, 2025).}

    \section{Premios, idiomas y software}
        \cvitem{2004}{Beca del Banco de la República para estudios de doctorado en el exterior.}
        \cvitem{2000}{Joven investigador, COLCIENCIAS.}
        \cvitem{Idiomas}{Español (lengua materna), inglés (C1).}
        \cvitem{Software}{Stata, R, EViews, Power BI, \LaTeX.}

\end{document}
